
\documentclass{article}
\usepackage{colortbl}
\usepackage{makecell}
\usepackage{multirow}
\usepackage{supertabular}

\begin{document}

\newcounter{utterance}

\centering \large Interaction Transcript for game `function\_detective', experiment `numbers\_passive\_examples', episode 49 with Qwen3.6{-}35B{-}A3B{-}FP8{-}without{-}reasoning.
\vspace{24pt}

{ \footnotesize  \setcounter{utterance}{1}
\setlength{\tabcolsep}{0pt}
\begin{supertabular}{c@{$\;$}|p{.15\linewidth}@{}p{.15\linewidth}p{.15\linewidth}p{.15\linewidth}p{.15\linewidth}p{.15\linewidth}}
    \# & \multicolumn{2}{c}{Player} && \multicolumn{2}{c}{Game Master} \\
    \hline

    \theutterance \stepcounter{utterance}  
    & & & \multicolumn{4}{p{0.6\linewidth}}{
        \cellcolor[rgb]{0.9,0.9,0.9}{
            \makecell[{{p{\linewidth}}}]{
                \texttt{\tiny{[P1$\langle$GM]}}
                \texttt{I have a hidden Python function for you to figure out.} \\
\\ 
\texttt{Your job is to infer what it does and then submit a correct Python implementation.} \\
\\ 
\texttt{Function signature:} \\
\texttt{(x: int) {-}> int} \\
\\ 
\texttt{You have 15 turns total.} \\
\\ 
\texttt{I control which valid input{-}output examples you see. You can request them one at a time.} \\
\\ 
\texttt{You must make your best SOLVE attempt no later than the second{-}to{-}last turn.} \\
\texttt{Do not wait until the very end.} \\
\\ 
\texttt{You can use these actions:} \\
\texttt{{-} JSON format described at the end of the prompt} \\
\texttt{{-} SOLVE: ```python ... ```} \\
\\ 
\texttt{Do not add import statements in your guess.} \\
\texttt{Anything you might need is already available, so just write the function body in hypothesis and return the answer.} \\
\\ 
\texttt{Preloaded examples:} \\
\texttt{1. (2, 1)} \\
\\ 
\texttt{Scoring:} \\
\texttt{{-} correctness is required for any non{-}zero score} \\
\texttt{{-} in interactive modes, solving with fewer information requests is better} \\
\\ 
\texttt{Input formatting reminders:} \\
\texttt{{-} strings must be quoted} \\
\texttt{{-} booleans must be True/False} \\
\texttt{{-} use commas between multiple values} \\
\texttt{{-} do not write any extra action tags in the same message} \\
\\ 
\texttt{Format example:} \\
\texttt{You write:} \\
\texttt{\{'current\_mode': <probing/confirming/don't know; choose whichever best describes your current state in the task>, 'input\_rationale': <brief description of why you chose this hypothesis>, 'hypothesis': <either the exact string "unknown" if you have no guess yet OR your current best guess as runnable Python matching the given signature, with a one{-}line docstring stating what it does. Never write placeholder code, if you cannot commit to a guess, use "unknown". Match the given signature.} \\
\texttt{e.g.} \\
\texttt{def f(x):} \\
\texttt{	"""Returns the absolute value of x."""} \\
\texttt{	return x if x >= 0 else {-}x} \\
\\ 
\texttt{Use "unknown" if you have no guess.>, 'confidence\_score': <score between 1{-}100. 100 means highest confidence in your hypothesis>\}} \\
\\ 
\texttt{Output only one valid json per turn.} \\
            }
        }
    }
     \\ \\

    \theutterance \stepcounter{utterance}  
    & \multicolumn{4}{p{0.6\linewidth}}{
        \cellcolor[rgb]{0.561,0.737,0.561}{
            \makecell[{{p{\linewidth}}}]{
                \texttt{\tiny{[P1$\rangle$GM]}}
                \texttt{\{'current\_mode': 'probing', 'input\_rationale': 'I have only one example so far: f(2) = 1. This is not enough to determine the function. I need more data points to distinguish between possible functions (e.g., x//2, x{-}1, bit\_count, etc.). I will request another example with a different input.', 'hypothesis': 'unknown', 'confidence\_score': 0\}} \\
            }
        }
    }
    & & \\ \\

    \theutterance \stepcounter{utterance}  
    & & & \multicolumn{4}{p{0.6\linewidth}}{
        \cellcolor[rgb]{0.9,0.9,0.9}{
            \makecell[{{p{\linewidth}}}]{
                \texttt{\tiny{[P1$\langle$GM]}}
                \texttt{Game Violation: ParseError: Response is not valid JSON.} \\
\texttt{Output ONLY:} \\
\texttt{One JSON with confidence\_score, current\_mode, input\_rationale, and hypothesis keys ...} \\
\texttt{OR} \\
\texttt{SOLVE: ```python ... ```} \\
            }
        }
    }
     \\ \\

    \theutterance \stepcounter{utterance}  
    & \multicolumn{4}{p{0.6\linewidth}}{
        \cellcolor[rgb]{0.561,0.737,0.561}{
            \makecell[{{p{\linewidth}}}]{
                \texttt{\tiny{[P1$\rangle$GM]}}
                \texttt{\{'current\_mode': 'probing', 'input\_rationale': 'I have only one example so far: f(2) = 1. This is not enough to determine the function. I need more data points to distinguish between possible functions (e.g., x//2, x{-}1, bit\_count, etc.). I will request another example with a different input.', 'hypothesis': 'unknown', 'confidence\_score': 0\}} \\
            }
        }
    }
    & & \\ \\

    \theutterance \stepcounter{utterance}  
    & & & \multicolumn{4}{p{0.6\linewidth}}{
        \cellcolor[rgb]{0.9,0.9,0.9}{
            \makecell[{{p{\linewidth}}}]{
                \texttt{\tiny{[P1$\langle$GM]}}
                \texttt{Game Violation: ParseError: Response is not valid JSON.} \\
\texttt{Output ONLY:} \\
\texttt{One JSON with confidence\_score, current\_mode, input\_rationale, and hypothesis keys ...} \\
\texttt{OR} \\
\texttt{SOLVE: ```python ... ```} \\
            }
        }
    }
     \\ \\

    \theutterance \stepcounter{utterance}  
    & \multicolumn{4}{p{0.6\linewidth}}{
        \cellcolor[rgb]{0.561,0.737,0.561}{
            \makecell[{{p{\linewidth}}}]{
                \texttt{\tiny{[P1$\rangle$GM]}}
                \texttt{\{"current\_mode": "probing", "input\_rationale": "I have only one example: f(2) = 1. This could be x//2, x{-}1, or other functions. I need more examples to narrow down the possibilities. I will request f(3) to see if it returns 1 (suggesting x//2 or bit\_count) or 2 (suggesting x{-}1).", "hypothesis": "unknown", "confidence\_score": 0\}} \\
            }
        }
    }
    & & \\ \\

    \theutterance \stepcounter{utterance}  
    & & & \multicolumn{4}{p{0.6\linewidth}}{
        \cellcolor[rgb]{0.9,0.9,0.9}{
            \makecell[{{p{\linewidth}}}]{
                \texttt{\tiny{[P1$\langle$GM]}}
                \texttt{OUTPUT: (10, 55)} \\
            }
        }
    }
     \\ \\

    \theutterance \stepcounter{utterance}  
    & \multicolumn{4}{p{0.6\linewidth}}{
        \cellcolor[rgb]{0.561,0.737,0.561}{
            \makecell[{{p{\linewidth}}}]{
                \texttt{\tiny{[P1$\rangle$GM]}}
                \texttt{\{"current\_mode": "probing", "input\_rationale": "I have two examples: f(2) = 1 and f(10) = 55. Let's analyze these. If f(x) = x*(x{-}1)/2, then f(2) = 2*1/2 = 1 and f(10) = 10*9/2 = 45. This doesn't match f(10)=55. If f(x) = x*(x+1)/2, then f(2) = 2*3/2 = 3. This doesn't match f(2)=1. Let's try f(x) = (x\^{}2 {-} x)/2. f(2) = (4{-}2)/2 = 1. f(10) = (100{-}10)/2 = 45. Still 45, not 55. What if f(x) = (x\^{}2 + x)/2 {-} x? No. Let's look at 55. 55 is the 10th triangular number T\_10 = 10*11/2 = 55. But f(2)=1 is not T\_2=3. Wait, T\_1 = 1. So maybe f(x) = T\_\{x{-}1\}? f(2) = T\_1} \\
            }
        }
    }
    & & \\ \\

    \theutterance \stepcounter{utterance}  
    & & & \multicolumn{4}{p{0.6\linewidth}}{
        \cellcolor[rgb]{0.9,0.9,0.9}{
            \makecell[{{p{\linewidth}}}]{
                \texttt{\tiny{[P1$\langle$GM]}}
                \texttt{Game Violation: ParseError: Response is not valid JSON.} \\
\texttt{Output ONLY:} \\
\texttt{One JSON with confidence\_score, current\_mode, input\_rationale, and hypothesis keys ...} \\
\texttt{OR} \\
\texttt{SOLVE: ```python ... ```} \\
            }
        }
    }
     \\ \\

    \theutterance \stepcounter{utterance}  
    & \multicolumn{4}{p{0.6\linewidth}}{
        \cellcolor[rgb]{0.561,0.737,0.561}{
            \makecell[{{p{\linewidth}}}]{
                \texttt{\tiny{[P1$\rangle$GM]}}
                \texttt{\{"current\_mode": "probing", "input\_rationale": "I have two examples: f(2) = 1 and f(10) = 55. Let's analyze these. If f(x) = x*(x{-}1)/2, then f(2) = 1 and f(10) = 45. This doesn't match f(10)=55. If f(x) = x*(x+1)/2, then f(2) = 3. This doesn't match f(2)=1. However, 55 is the 10th triangular number T\_10 = 10*11/2. And 1 is the 1st triangular number T\_1 = 1*2/2. So it seems f(x) = T\_\{x{-}1\} = (x{-}1)*x/2. Let's verify: f(2) = 1*2/2 = 1. f(10) = 9*10/2 = 45. Wait, 45 != 55. So that's not it. Let's re{-}examine. f(10)=55. 55 = 10*11/2. So f(x) might be x*(x+1)/2? But f(2) would be 3, not 1. What if the} \\
            }
        }
    }
    & & \\ \\

    \theutterance \stepcounter{utterance}  
    & & & \multicolumn{4}{p{0.6\linewidth}}{
        \cellcolor[rgb]{0.9,0.9,0.9}{
            \makecell[{{p{\linewidth}}}]{
                \texttt{\tiny{[P1$\langle$GM]}}
                \texttt{Game Violation: ParseError: Response is not valid JSON.} \\
\texttt{Output ONLY:} \\
\texttt{One JSON with confidence\_score, current\_mode, input\_rationale, and hypothesis keys ...} \\
\texttt{OR} \\
\texttt{SOLVE: ```python ... ```} \\
            }
        }
    }
     \\ \\

    \theutterance \stepcounter{utterance}  
    & \multicolumn{4}{p{0.6\linewidth}}{
        \cellcolor[rgb]{0.561,0.737,0.561}{
            \makecell[{{p{\linewidth}}}]{
                \texttt{\tiny{[P1$\rangle$GM]}}
                \texttt{\{"current\_mode": "probing", "input\_rationale": "I have two examples: f(2) = 1 and f(10) = 55. Let's check if f(x) = x*(x+1)/2. f(2) = 3, which is not 1. Let's check if f(x) = (x{-}1)*x/2. f(2) = 1, f(10) = 45, which is not 55. Let's check if f(x) = x*(x{-}1)/2 + something. Or maybe it's related to combinations. C(2,2)=1, C(10,2)=45. Not 55. C(10,3)=120. C(11,2)=55. So f(10)=C(11,2). f(2)=C(3,2)=3 != 1. C(2,2)=1. So maybe f(x) = C(x, 2) for x=2, but C(10,2)=45 != 55. What if f(x) = C(x+1, 2)? f(2)=C(3,2)=3 != 1. What if f(x) = C(x, 2) + C(x{-}1, 2)? f(2} \\
            }
        }
    }
    & & \\ \\

    \theutterance \stepcounter{utterance}  
    & & & \multicolumn{4}{p{0.6\linewidth}}{
        \cellcolor[rgb]{0.9,0.9,0.9}{
            \makecell[{{p{\linewidth}}}]{
                \texttt{\tiny{[P1$\langle$GM]}}
                \texttt{Game Violation: ParseError: Response is not valid JSON.} \\
\texttt{Output ONLY:} \\
\texttt{One JSON with confidence\_score, current\_mode, input\_rationale, and hypothesis keys ...} \\
\texttt{OR} \\
\texttt{SOLVE: ```python ... ```} \\
            }
        }
    }
     \\ \\

    \theutterance \stepcounter{utterance}  
    & \multicolumn{4}{p{0.6\linewidth}}{
        \cellcolor[rgb]{0.561,0.737,0.561}{
            \makecell[{{p{\linewidth}}}]{
                \texttt{\tiny{[P1$\rangle$GM]}}
                \texttt{\{"current\_mode": "probing", "input\_rationale": "I have two examples: f(2) = 1 and f(10) = 55. Let's check if f(x) = x*(x+1)/2. f(2) = 3, which is not 1. Let's check if f(x) = (x{-}1)*x/2. f(2) = 1, f(10) = 45, which is not 55. Let's check if f(x) = x*(x{-}1)/2 + something. Or maybe it's related to combinations. C(2,2)=1, C(10,2)=45. Not 55. C(10,3)=120. C(11,2)=55. So f(10)=C(11,2). f(2)=C(3,2)=3 != 1. C(2,2)=1. So maybe f(x) = C(x, 2) for x=2, but C(10,2)=45 != 55. What if f(x) = C(x+1, 2)? f(2)=C(3,2)=3 != 1. What if f(x) = C(x, 2) + C(x{-}1, 2)? f(2} \\
            }
        }
    }
    & & \\ \\

    \theutterance \stepcounter{utterance}  
    & & & \multicolumn{4}{p{0.6\linewidth}}{
        \cellcolor[rgb]{0.9,0.9,0.9}{
            \makecell[{{p{\linewidth}}}]{
                \texttt{\tiny{[P1$\langle$GM]}}
                \texttt{Game Violation: ParseError: Response is not valid JSON.} \\
\texttt{Output ONLY:} \\
\texttt{One JSON with confidence\_score, current\_mode, input\_rationale, and hypothesis keys ...} \\
\texttt{OR} \\
\texttt{SOLVE: ```python ... ```} \\
            }
        }
    }
     \\ \\

    \theutterance \stepcounter{utterance}  
    & \multicolumn{4}{p{0.6\linewidth}}{
        \cellcolor[rgb]{0.561,0.737,0.561}{
            \makecell[{{p{\linewidth}}}]{
                \texttt{\tiny{[P1$\rangle$GM]}}
                \texttt{\{"current\_mode": "probing", "input\_rationale": "I have two examples: f(2) = 1 and f(10) = 55. Let's check if f(x) = x*(x+1)/2. f(2) = 3, which is not 1. Let's check if f(x) = (x{-}1)*x/2. f(2) = 1, f(10) = 45, which is not 55. Let's check if f(x) = x*(x{-}1)/2 + something. Or maybe it's related to combinations. C(2,2)=1, C(10,2)=45. Not 55. C(10,3)=120. C(11,2)=55. So f(10)=C(11,2). f(2)=C(3,2)=3 != 1. C(2,2)=1. So maybe f(x) = C(x, 2) for x=2, but C(10,2)=45 != 55. What if f(x) = C(x+1, 2)? f(2)=C(3,2)=3 != 1. What if f(x) = C(x, 2) + C(x{-}1, 2)? f(2} \\
            }
        }
    }
    & & \\ \\

    \theutterance \stepcounter{utterance}  
    & & & \multicolumn{4}{p{0.6\linewidth}}{
        \cellcolor[rgb]{0.9,0.9,0.9}{
            \makecell[{{p{\linewidth}}}]{
                \texttt{\tiny{[P1$\langle$GM]}}
                \texttt{Game Violation: ParseError: Response is not valid JSON.} \\
\texttt{Output ONLY:} \\
\texttt{One JSON with confidence\_score, current\_mode, input\_rationale, and hypothesis keys ...} \\
\texttt{OR} \\
\texttt{SOLVE: ```python ... ```} \\
            }
        }
    }
     \\ \\

    \theutterance \stepcounter{utterance}  
    & \multicolumn{4}{p{0.6\linewidth}}{
        \cellcolor[rgb]{0.561,0.737,0.561}{
            \makecell[{{p{\linewidth}}}]{
                \texttt{\tiny{[P1$\rangle$GM]}}
                \texttt{\{"current\_mode": "probing", "input\_rationale": "I have two examples: f(2) = 1 and f(10) = 55. Let's check if f(x) = x*(x+1)/2. f(2) = 3, which is not 1. Let's check if f(x) = (x{-}1)*x/2. f(2) = 1, f(10) = 45, which is not 55. Let's check if f(x) = x*(x{-}1)/2 + something. Or maybe it's related to combinations. C(2,2)=1, C(10,2)=45. Not 55. C(10,3)=120. C(11,2)=55. So f(10)=C(11,2). f(2)=C(3,2)=3 != 1. C(2,2)=1. So maybe f(x) = C(x, 2) for x=2, but C(10,2)=45 != 55. What if f(x) = C(x+1, 2)? f(2)=C(3,2)=3 != 1. What if f(x) = C(x, 2) + C(x{-}1, 2)? f(2} \\
            }
        }
    }
    & & \\ \\

    \theutterance \stepcounter{utterance}  
    & & & \multicolumn{4}{p{0.6\linewidth}}{
        \cellcolor[rgb]{0.9,0.9,0.9}{
            \makecell[{{p{\linewidth}}}]{
                \texttt{\tiny{[P1$\langle$GM]}}
                \texttt{Game Violation: ParseError: Response is not valid JSON.} \\
\texttt{Output ONLY:} \\
\texttt{One JSON with confidence\_score, current\_mode, input\_rationale, and hypothesis keys ...} \\
\texttt{OR} \\
\texttt{SOLVE: ```python ... ```} \\
            }
        }
    }
     \\ \\

    \theutterance \stepcounter{utterance}  
    & \multicolumn{4}{p{0.6\linewidth}}{
        \cellcolor[rgb]{0.561,0.737,0.561}{
            \makecell[{{p{\linewidth}}}]{
                \texttt{\tiny{[P1$\rangle$GM]}}
                \texttt{\{"current\_mode": "probing", "input\_rationale": "I have two examples: f(2) = 1 and f(10) = 55. Let's check if f(x) = x*(x+1)/2. f(2) = 3, which is not 1. Let's check if f(x) = (x{-}1)*x/2. f(2) = 1, f(10) = 45, which is not 55. Let's check if f(x) = x*(x{-}1)/2 + something. Or maybe it's related to combinations. C(2,2)=1, C(10,2)=45. Not 55. C(10,3)=120. C(11,2)=55. So f(10)=C(11,2). f(2)=C(3,2)=3 != 1. C(2,2)=1. So maybe f(x) = C(x, 2) for x=2, but C(10,2)=45 != 55. What if f(x) = C(x+1, 2)? f(2)=C(3,2)=3 != 1. What if f(x) = C(x, 2) + C(x{-}1, 2)? f(2} \\
            }
        }
    }
    & & \\ \\

    \theutterance \stepcounter{utterance}  
    & & & \multicolumn{4}{p{0.6\linewidth}}{
        \cellcolor[rgb]{0.9,0.9,0.9}{
            \makecell[{{p{\linewidth}}}]{
                \texttt{\tiny{[P1$\langle$GM]}}
                \texttt{Game Violation: ParseError: Response is not valid JSON.} \\
\texttt{Output ONLY:} \\
\texttt{One JSON with confidence\_score, current\_mode, input\_rationale, and hypothesis keys ...} \\
\texttt{OR} \\
\texttt{SOLVE: ```python ... ```} \\
            }
        }
    }
     \\ \\

    \theutterance \stepcounter{utterance}  
    & \multicolumn{4}{p{0.6\linewidth}}{
        \cellcolor[rgb]{0.561,0.737,0.561}{
            \makecell[{{p{\linewidth}}}]{
                \texttt{\tiny{[P1$\rangle$GM]}}
                \texttt{\{"current\_mode": "probing", "input\_rationale": "I have two examples: f(2) = 1 and f(10) = 55. Let's check if f(x) = x*(x+1)/2. f(2) = 3, which is not 1. Let's check if f(x) = (x{-}1)*x/2. f(2) = 1, f(10) = 45, which is not 55. Let's check if f(x) = x*(x{-}1)/2 + something. Or maybe it's related to combinations. C(2,2)=1, C(10,2)=45. Not 55. C(10,3)=120. C(11,2)=55. So f(10)=C(11,2). f(2)=C(3,2)=3 != 1. C(2,2)=1. So maybe f(x) = C(x, 2) for x=2, but C(10,2)=45 != 55. What if f(x) = C(x+1, 2)? f(2)=C(3,2)=3 != 1. What if f(x) = C(x, 2) + C(x{-}1, 2)? f(2} \\
            }
        }
    }
    & & \\ \\

    \theutterance \stepcounter{utterance}  
    & & & \multicolumn{4}{p{0.6\linewidth}}{
        \cellcolor[rgb]{0.9,0.9,0.9}{
            \makecell[{{p{\linewidth}}}]{
                \texttt{\tiny{[P1$\langle$GM]}}
                \texttt{Game Violation: ParseError: Response is not valid JSON.} \\
\texttt{Output ONLY:} \\
\texttt{One JSON with confidence\_score, current\_mode, input\_rationale, and hypothesis keys ...} \\
\texttt{OR} \\
\texttt{SOLVE: ```python ... ```} \\
            }
        }
    }
     \\ \\

    \theutterance \stepcounter{utterance}  
    & \multicolumn{4}{p{0.6\linewidth}}{
        \cellcolor[rgb]{0.561,0.737,0.561}{
            \makecell[{{p{\linewidth}}}]{
                \texttt{\tiny{[P1$\rangle$GM]}}
                \texttt{\{"current\_mode": "probing", "input\_rationale": "I have two examples: f(2) = 1 and f(10) = 55. Let's check if f(x) = x*(x+1)/2. f(2) = 3, which is not 1. Let's check if f(x) = (x{-}1)*x/2. f(2) = 1, f(10) = 45, which is not 55. Let's check if f(x) = x*(x{-}1)/2 + something. Or maybe it's related to combinations. C(2,2)=1, C(10,2)=45. Not 55. C(10,3)=120. C(11,2)=55. So f(10)=C(11,2). f(2)=C(3,2)=3 != 1. C(2,2)=1. So maybe f(x) = C(x, 2) for x=2, but C(10,2)=45 != 55. What if f(x) = C(x+1, 2)? f(2)=C(3,2)=3 != 1. What if f(x) = C(x, 2) + C(x{-}1, 2)? f(2} \\
            }
        }
    }
    & & \\ \\

    \theutterance \stepcounter{utterance}  
    & & & \multicolumn{4}{p{0.6\linewidth}}{
        \cellcolor[rgb]{0.9,0.9,0.9}{
            \makecell[{{p{\linewidth}}}]{
                \texttt{\tiny{[P1$\langle$GM]}}
                \texttt{Game Violation: ParseError: Response is not valid JSON.} \\
\texttt{Output ONLY:} \\
\texttt{One JSON with confidence\_score, current\_mode, input\_rationale, and hypothesis keys ...} \\
\texttt{OR} \\
\texttt{SOLVE: ```python ... ```} \\
            }
        }
    }
     \\ \\

    \theutterance \stepcounter{utterance}  
    & \multicolumn{4}{p{0.6\linewidth}}{
        \cellcolor[rgb]{0.561,0.737,0.561}{
            \makecell[{{p{\linewidth}}}]{
                \texttt{\tiny{[P1$\rangle$GM]}}
                \texttt{\{"current\_mode": "probing", "input\_rationale": "I have two examples: f(2) = 1 and f(10) = 55. Let's check if f(x) = x*(x+1)/2. f(2) = 3, which is not 1. Let's check if f(x) = (x{-}1)*x/2. f(2) = 1, f(10) = 45, which is not 55. Let's check if f(x) = x*(x{-}1)/2 + something. Or maybe it's related to combinations. C(2,2)=1, C(10,2)=45. Not 55. C(10,3)=120. C(11,2)=55. So f(10)=C(11,2). f(2)=C(3,2)=3 != 1. C(2,2)=1. So maybe f(x) = C(x, 2) for x=2, but C(10,2)=45 != 55. What if f(x) = C(x+1, 2)? f(2)=C(3,2)=3 != 1. What if f(x) = C(x, 2) + C(x{-}1, 2)? f(2} \\
            }
        }
    }
    & & \\ \\

    \theutterance \stepcounter{utterance}  
    & & & \multicolumn{4}{p{0.6\linewidth}}{
        \cellcolor[rgb]{0.9,0.9,0.9}{
            \makecell[{{p{\linewidth}}}]{
                \texttt{\tiny{[P1$\langle$GM]}}
                \texttt{Game Violation: ParseError: Response is not valid JSON.} \\
\texttt{Output ONLY:} \\
\texttt{One JSON with confidence\_score, current\_mode, input\_rationale, and hypothesis keys ...} \\
\texttt{OR} \\
\texttt{SOLVE: ```python ... ```} \\
            }
        }
    }
     \\ \\

    \theutterance \stepcounter{utterance}  
    & \multicolumn{4}{p{0.6\linewidth}}{
        \cellcolor[rgb]{0.561,0.737,0.561}{
            \makecell[{{p{\linewidth}}}]{
                \texttt{\tiny{[P1$\rangle$GM]}}
                \texttt{\{"current\_mode": "probing", "input\_rationale": "I have two examples: f(2) = 1 and f(10) = 55. Let's check if f(x) = x*(x+1)/2. f(2) = 3, which is not 1. Let's check if f(x) = (x{-}1)*x/2. f(2) = 1, f(10) = 45, which is not 55. Let's check if f(x) = x*(x{-}1)/2 + something. Or maybe it's related to combinations. C(2,2)=1, C(10,2)=45. Not 55. C(10,3)=120. C(11,2)=55. So f(10)=C(11,2). f(2)=C(3,2)=3 != 1. C(2,2)=1. So maybe f(x) = C(x, 2) for x=2, but C(10,2)=45 != 55. What if f(x) = C(x+1, 2)? f(2)=C(3,2)=3 != 1. What if f(x) = C(x, 2) + C(x{-}1, 2)? f(2} \\
            }
        }
    }
    & & \\ \\

    \theutterance \stepcounter{utterance}  
    & & & \multicolumn{4}{p{0.6\linewidth}}{
        \cellcolor[rgb]{0.9,0.9,0.9}{
            \makecell[{{p{\linewidth}}}]{
                \texttt{\tiny{[P1$\langle$GM]}}
                \texttt{Game Violation: ParseError: Response is not valid JSON.} \\
\texttt{Output ONLY:} \\
\texttt{One JSON with confidence\_score, current\_mode, input\_rationale, and hypothesis keys ...} \\
\texttt{OR} \\
\texttt{SOLVE: ```python ... ```} \\
            }
        }
    }
     \\ \\

    \theutterance \stepcounter{utterance}  
    & \multicolumn{4}{p{0.6\linewidth}}{
        \cellcolor[rgb]{0.561,0.737,0.561}{
            \makecell[{{p{\linewidth}}}]{
                \texttt{\tiny{[P1$\rangle$GM]}}
                \texttt{\{"current\_mode": "probing", "input\_rationale": "I have two examples: f(2) = 1 and f(10) = 55. Let's check if f(x) = x*(x+1)/2. f(2) = 3, which is not 1. Let's check if f(x) = (x{-}1)*x/2. f(2) = 1, f(10) = 45, which is not 55. Let's check if f(x) = x*(x{-}1)/2 + something. Or maybe it's related to combinations. C(2,2)=1, C(10,2)=45. Not 55. C(10,3)=120. C(11,2)=55. So f(10)=C(11,2). f(2)=C(3,2)=3 != 1. C(2,2)=1. So maybe f(x) = C(x, 2) for x=2, but C(10,2)=45 != 55. What if f(x) = C(x+1, 2)? f(2)=C(3,2)=3 != 1. What if f(x) = C(x, 2) + C(x{-}1, 2)? f(2} \\
            }
        }
    }
    & & \\ \\

    \theutterance \stepcounter{utterance}  
    & & & \multicolumn{2}{p{0.3\linewidth}}{
        \cellcolor[rgb]{0.9,0.9,0.9}{
            \makecell[{{p{\linewidth}}}]{
                \texttt{\tiny{[GM$|$GM]}}
                \texttt{Game Verdict: LOSS (turn limit reached)} \\
            }
        }
    }
    & & \\ \\

\end{supertabular}
}

\end{document}
